\chapter{System Architecture}

\section{Architectural objectives and invariants}

The system is designed for a regional research and education setting in which a small team must
inspect EO metadata, reproduce experiments, and remain able to replace an external provider
without rewriting the domain logic. Five invariants make that design operational rather than
merely aspirational.

\begin{description}
\item[Boundary validation] Untrusted provider JSON is validated before transformation or
persistence. Every stored scene therefore satisfies the same identity, temporal, geometry, asset,
and quality-shape contracts, irrespective of its source.
\item[Transport decoupling] HTTP, OAuth2, fixture, and test transports implement injectable
boundaries. Search, transformation, catalog, and report code do not depend on a particular
network client or provider SDK.
\item[Bounded concurrency] The number of active ingestion workers is explicit and enforced by an
\texttt{asyncio.Semaphore}. If the configured limit is \(k\), the invariant
\(0\leq A(t)\leq k\) holds at every instant \(t\), protecting both provider rate limits and local
resources.
\item[Append-only provenance] Raw payloads are content-addressed and manifest records are appended
rather than rewritten. An observation can therefore be audited without trusting mutable catalog
metadata.
\item[Deterministic spatial evaluation] Coordinate assumptions, bounding-box prefilters, exact
polygon checks, and boundary semantics are explicit. Identical metadata produces identical
regional inclusion results.
\end{description}

These invariants form a chain: validation makes persistence safe, transport decoupling makes
validation testable, bounded concurrency makes provider interaction predictable, provenance makes
the result explainable, and deterministic spatial evaluation makes regional claims reproducible.
Breaking one invariant weakens the interpretation of all downstream metrics.

\section{Layered system decomposition}

The architecture separates six responsibilities:

\begin{enumerate}
\item \textbf{Provider adapters} issue STAC searches, follow pagination, handle OAuth2, and
translate transient transport failures into bounded retry decisions.
\item \textbf{Boundary models} validate and normalize STAC Items into typed \texttt{EOScene}
records with a stable shape independent of provider-specific JSON extensions.
\item \textbf{Orchestration} schedules transformations and persistence under a worker limit while
collecting elapsed time, payload, retry, and failure metrics.
\item \textbf{Evidence storage} writes raw canonical payloads and JSON Lines provenance, then stores
queryable scene metadata in the SQLite \texttt{SpatialCatalog}.
\item \textbf{Analytical services} profile metadata quality, filter by cloud cover and required
assets, evaluate regional geometry, fetch byte ranges, and calculate NDVI.
\item \textbf{Interfaces and reporting} expose read-only educational HTTP queries and produce
deterministic JSON/Markdown reports, including additive dated live-provider artifacts.
\end{enumerate}

The layers are connected by typed values and narrow protocols rather than shared global state.
Consequently, an offline fixture can exercise the same transformation and persistence path as a
live CDSE response.

\section{End-to-end data transformation}

The principal flow is shown below. The arrows denote ownership boundaries, not merely function
calls.

\begin{verbatim}
 STAC/CDSE endpoint or offline fixture
              |
              v
 injectable transport -- OAuth2 / RetryPolicy / pagination
              |
              v
 validated STACItem -- schema boundary
              |
              v
 transform_stac_to_scene --> typed EOScene
              |                         |
              |                         +--> quality profile
              v
 SHA-256 raw payload + JSONL manifest
              |
              v
 ConcurrentEOIngestor -- bounded tasks --> SQLite SpatialCatalog
                                              |             |
                                              v             v
                                      ROI/quality queries  reports
                                                            |
                                                            v
                                                    educational HTTP API
\end{verbatim}

At the provider boundary, a search response is paginated and each Item is validated independently.
The transformation preserves the identifier, acquisition time, platform, geometry, bounding box,
EPSG code, cloud metadata, and asset links. Persistence then has two complementary products:
the provenance store preserves what arrived, while the catalog stores the indexed metadata needed
for low-latency query. Reports consume metrics from the orchestration path and counts from the
catalog, making inconsistencies observable.

\section{Provider and ingestion contracts}

\texttt{STACClient} owns request construction, pagination links, response-shape checks, timeout
classification, and 429 handling. \texttt{CDSEClient} composes it with
\texttt{CDSETokenProvider}; credentials are resolved from environment variables and token refresh
is protected by an asynchronous lock. Invalid schemas are permanent errors and are not retried.
Timeouts and rate limits use a bounded exponential policy, optionally honoring a provider
\texttt{Retry-After} value.

\texttt{ConcurrentEOIngestor} owns task scheduling, not provider policy. It acquires its semaphore
before expensive work and releases it in a \texttt{finally} path. Blocking filesystem and SQLite
operations are delegated to worker threads, while provenance appends are serialized by an
asynchronous lock. This division prevents an individual slow write from blocking unrelated
network tasks without allowing writes to race on the manifest.

\section{Kvarken regional filtering contract}

The regional preset is an EPSG:4326 WGS 84 bounding box
\[
 B_K=[20.5,22.5]\times[62.8,63.8],
\]
where the first coordinate is longitude in degrees east and the second is latitude in degrees
north. The order is part of the contract; coordinates in another CRS are rejected rather than
silently reprojected.

For scene rectangle
\[
B_S=[x^-_S,x^+_S]\times[y^-_S,y^+_S],
\]
the catalog prefilter accepts a candidate exactly when
\[
x^-_S\leq x^+_K\ \land\ x^+_S\geq x^-_K\ \land\
y^-_S\leq y^+_K\ \land\ y^+_S\geq y^-_K.
\]
The inequalities are non-strict so contact at a boundary is retained. SQLite applies this
predicate over indexed columns, reducing the candidate set before exact polygon processing.

The exact contract is polygonal intersection:
\[
\operatorname{in}(S,K)=
\begin{cases}
1,&\operatorname{intersects}(\operatorname{footprint}(S),K),\\
0,&\text{otherwise}.
\end{cases}
\]
Point-in-polygon uses ray casting with an odd crossing count, and a point on an edge is treated as
intersecting. This two-stage approach combines predictable SQL performance with a deterministic
geometry result, while avoiding a false claim that rectangular overlap alone proves polygon
overlap.

\section{Persistence, API, and educational separation}

\texttt{SpatialCatalog} is a WAL-enabled SQLite store with indexes over identity, platform,
acquisition time, cloud cover, and bounding-box coordinates. It stores normalized metadata and
serialized geometry/assets, not the raw provider document. The raw document remains in the
content-addressed provenance area, so catalog maintenance, retention, or a teaching query cannot
silently rewrite the evidence used to generate a report.

The educational HTTP layer is intentionally read-only. Health checks verify that the local
pipeline is operational, while scene endpoints expose bounded JSON query results over catalog
metadata. It does not expose credentials, raw payload paths, or arbitrary SQL. This separation
keeps a simple browser-facing interface useful for demonstrations without turning the provenance
store into an application database.

The raster slice follows the same boundary principle: it requests inclusive HTTP byte ranges and
calculates NDVI over supplied aligned values, but does not pretend to be a full GeoTIFF server.
The result is a compact architecture in which research reproducibility and educational
accessibility reinforce rather than compromise operational safety.

\section{Operational failure boundaries}

Failures are classified at the narrowest layer that can interpret them. A timeout or HTTP 429 is a
transient provider event; bounded retry may recover it. A malformed JSON response is a permanent
input or protocol failure; it is surfaced with context. A SQLite transaction error rolls back the
write, while an already committed provenance object remains inspectable. A missing CDSE
configuration is not an error in offline mode: the runner records fixture mode explicitly and
preserves the same report schema. These choices ensure that the system fails visibly without
making CI or teaching depend on external credentials.
